\begin{table}[ht]
\centering
\footnotesize
\caption{逐格校准分类计数（每行 $n=28$）。完整 112 行拒绝率、Wilson 区间与分类结果见补充材料。}\label{tab:calibration}
\begin{tabular}{llcccc}
\toprule
表示 & 方法 & 实用校准 & 实质保守 & 实质反保守 & 边界/未定\\
\midrule
BPI & M0 & 7 (25.0) & 0 (0.0) & 20 (71.4) & 1 (3.6)\\
BPI & M1 & 5 (17.9) & 10 (35.7) & 7 (25.0) & 6 (21.4)\\
MOMENTS & M0 & 6 (21.4) & 0 (0.0) & 20 (71.4) & 2 (7.1)\\
MOMENTS & M1 & 1 (3.6) & 12 (42.9) & 7 (25.0) & 8 (28.6)\\
\bottomrule
\end{tabular}

\begin{minipage}{0.94\textwidth}
\vspace{0.4em}
\footnotesize 注：括号内为该类计数占 28 格的百分比。实用校准采用 90\% Wilson 区间完整落在等效带 $[0.035,0.065]$ 内，对应预设的实用等效性判据（两个单侧 5\% 判定等价于一个 90\% 双侧区间）；实质保守要求 95\% 区间完整位于实用带下方，实质反保守要求 95\% 区间完整位于实用带上方；边界/未定为既不满足实用校准、也不满足实质保守或实质反保守的剩余状态。上述计数为预设设计网格内的描述性分布，不构成方法在更广参数空间中的通过率。
\end{minipage}
\end{table}
